% =====================================================================
% Stage 2: Methodology and Experimental Setup
% Required packages: amsmath, amssymb, bm, booktabs, siunitx, graphicx
% Notation: codes a, b, c in \mathcal{C} = {PSN, pRT, TRX}
%   PSN = POSEIDON, pRT = petitRADTRANS, TRX = TauREx
% =====================================================================

\section{Methodology}
\label{sec:method}

\subsection{Problem definition}
\label{sec:problem}

Let $\bm{\theta}\in\Theta\subset\mathbb{R}^{9}$ denote the parameters of a transiting planet,
\begin{equation}
  \bm{\theta} = \bigl(R_\star,\; R_p,\; M_p,\; T,\; \log_{10}X_{\mathrm{H_2O}},\; \log_{10}X_{\mathrm{CH_4}},\; \log_{10}X_{\mathrm{CO}},\; \log_{10}X_{\mathrm{CO_2}},\; \log_{10}X_{\mathrm{NH_3}}\bigr),
  \label{eq:theta}
\end{equation}
where $R_\star$ is the stellar radius, $R_p$ and $M_p$ are the planetary radius and mass, $T$ is the temperature of an isothermal atmosphere, and $X_s$ is the volume mixing ratio of species $s$.
A radiative-transfer code $c$ defines a forward model
\begin{equation}
  \mathcal{F}^{(c)}:\Theta\rightarrow\mathbb{R}^{D},\qquad \mathbf{y}^{(c)} = \mathcal{F}^{(c)}(\bm{\theta}),
  \label{eq:forward}
\end{equation}
which maps $\bm{\theta}$ to a transmission spectrum of fractional transit depths in $D=52$ wavelength bins.
We consider three codes, $c\in\mathcal{C}=\{\mathrm{PSN},\mathrm{pRT},\mathrm{TRX}\}$, for POSEIDON~\cite{TODO_macdonald_poseidon}, petitRADTRANS~\cite{TODO_molliere_petitradtrans}, and TauREx~\cite{TODO_alrefaie_taurex3}.
For a pair of codes $(a,b)$ and the same $\bm{\theta}$, the inter-code systematic is the difference spectrum
\begin{equation}
  \bm{\Delta}^{(ab)}(\bm{\theta}) = \mathcal{F}^{(a)}(\bm{\theta}) - \mathcal{F}^{(b)}(\bm{\theta}) \in \mathbb{R}^{D}.
  \label{eq:delta}
\end{equation}
In a physically consistent setting with identical inputs, $\bm{\Delta}^{(ab)}\equiv\mathbf{0}$.
We address three questions:
\begin{enumerate}
  \item[(Q1)] \emph{Structure.} What are the amplitude, spectral shape, and parameter dependence of $\bm{\Delta}^{(ab)}$?
  \item[(Q2)] \emph{Detectability.} Given a single spectrum, possibly corrupted by observational noise, can one determine which code produced it?
  \item[(Q3)] \emph{Impact on inference.} How much bias does a data-driven retrieval model trained on code $a$ incur when it is applied to spectra from code $b\neq a$?
\end{enumerate}
We treat (Q1) with a paired residual decomposition (Section~\ref{sec:residuals}), (Q2) with one-class autoencoders and a supervised classifier (Sections~\ref{sec:ae}--\ref{sec:rf}), and (Q3) with gradient-boosted regression (Section~\ref{sec:retrieval}).

\subsection{Simulated spectra}
\label{sec:data}

\paragraph{Parameter sampling.}
We draw $N=98{,}725$ parameter vectors [TODO: VERIFY whether 98,725 is the full sample or the sample remaining after a cleaning step, and what that step removed].
In the stored sample, $R_\star$, $R_p$, $M_p$, and $T$ are approximately uniform on $[0.5,2.0]\,R_\odot$, $[0.1,2.0]\,R_\mathrm{J}$, $[0.01,5.0]\,M_\mathrm{J}$ [TODO: VERIFY lower bound of $M_p$; the lowest-mass decile is under-populated by $\approx$12\%], and $[300,3000]\,\mathrm{K}$, and the five log-abundances are uniform on $[-9,-3]$.
Two abundances are restricted where opacity data are unavailable at high temperature: $\log_{10}X_{\mathrm{CH_4}}\in[-9,-7]$ for $T>2000\,\mathrm{K}$ and $\log_{10}X_{\mathrm{NH_3}}\in[-9,-7]$ for $T>1600\,\mathrm{K}$ [TODO: VERIFY; author notes give 1500\,K, but the stored parameters show the cut-offs at 2000\,K for CH$_4$ and 1600\,K for NH$_3$].
This restriction induces negative correlations between $T$ and the two log-abundances (Pearson $r=-0.47$ for CH$_4$ and $-0.53$ for NH$_3$) and a positive correlation between the two abundances themselves ($r=0.26$); all other parameter pairs are uncorrelated ($|r|<0.02$).

\paragraph{Forward models.}
All three codes use an isothermal pressure--temperature profile at $T$, a cloud-free atmosphere, and each code's default opacity database [TODO: VERIFY line lists, collision-induced absorption, Rayleigh scattering, background gas composition, and reference pressure for each code].
POSEIDON and TauREx compute spectra from cross-sections at a native resolving power of $R=15{,}000$; petitRADTRANS uses correlated-$k$ opacities at $R=1000$.
Each native spectrum is binned onto a common grid of $D=52$ bins that follows the Ariel spectral coverage~\cite{TODO_tinetti_ariel_mission} [TODO: VERIFY Ariel tier and binning reference].
Bin centres span $0.55$--$7.28\,\mu\mathrm{m}$: three photometric bins (0.55, 0.70, 0.95\,$\mu$m), followed by three spectroscopic segments with approximately constant $\lambda/\Delta\lambda\approx 11$ ($1.16$--$1.88\,\mu$m, 6 bins), $\approx 51$ ($1.97$--$3.72\,\mu$m, 33 bins), and $\approx 16$ ($4.03$--$7.28\,\mu$m, 10 bins).
Because the three codes share $\bm{\theta}$ row by row, every planet $i$ has a triplet $\bigl(\mathbf{y}_i^{(\mathrm{PSN})},\mathbf{y}_i^{(\mathrm{pRT})},\mathbf{y}_i^{(\mathrm{TRX})}\bigr)$, so the effect of the code can be isolated from that of the planet.

\paragraph{Feature-height selection.}
We define the feature height of a spectrum as
\begin{equation}
  h(\mathbf{y}) = \max_{j} y_j - \min_{j} y_j ,
  \label{eq:fh}
\end{equation}
and retain planet $i$ only if $h(\mathbf{y}_i^{(c)})\geq h_{\min}=210\,\mathrm{ppm}$ for every code under consideration [TODO: VERIFY physical justification for 210\,ppm].
The median feature height of the unselected sample is 153--156\,ppm, depending on the code, so the selection retains 44.6\% of planets.

\paragraph{Per-spectrum normalisation.}
For $\mathbf{y}\in\mathbb{R}^{D}$ we define the bin-averaged depth and the population standard deviation
\begin{equation}
  m(\mathbf{y}) = \frac{1}{D}\sum_{j=1}^{D} y_j,\qquad
  s(\mathbf{y}) = \Bigl(\frac{1}{D}\sum_{j=1}^{D}\bigl(y_j-m(\mathbf{y})\bigr)^2\Bigr)^{1/2},
  \label{eq:ms}
\end{equation}
and the standardised spectrum $\tilde{\mathbf{y}} = \bigl(\mathbf{y}-m(\mathbf{y})\mathbf{1}\bigr)/s(\mathbf{y})$.
Standardisation removes the absolute depth, which encodes $(R_p/R_\star)^2$, and the overall feature amplitude.
Models that operate on $\tilde{\mathbf{y}}$ are therefore sensitive only to spectral shape; where absolute depth is required (Section~\ref{sec:retrieval}), $m$ and $s$ are supplied as separate inputs.

\paragraph{Noise model.}
Observational noise is modelled as white Gaussian noise with a uniform standard deviation $\sigma$ in every bin,
\begin{equation}
  \mathbf{y}^{\sigma} = \mathbf{y} + \bm{\epsilon},\qquad \bm{\epsilon}\sim\mathcal{N}\bigl(\mathbf{0},\sigma^{2}\mathbf{I}_{D}\bigr),
  \label{eq:noise}
\end{equation}
added to the transit-depth spectrum before standardisation, so that $\tilde{\mathbf{y}}^{\sigma}$ is computed with $m(\mathbf{y}^{\sigma})$ and $s(\mathbf{y}^{\sigma})$.

\subsection{Paired residual decomposition (Q1)}
\label{sec:residuals}

For each planet $i$ and code pair $(a,b)$ we split $\bm{\Delta}_i^{(ab)}$ (Eq.~\ref{eq:delta}) into a wavelength-independent offset and a shape residual,
\begin{equation}
  \delta_i^{(ab)} = \frac{1}{D}\mathbf{1}^{\top}\bm{\Delta}_i^{(ab)},\qquad
  \mathbf{r}_i^{(ab)} = \bm{\Delta}_i^{(ab)} - \delta_i^{(ab)}\mathbf{1}.
  \label{eq:offset}
\end{equation}
The offset is degenerate with $R_p$ in a retrieval and is reported separately.
We summarise the shape residual by its root-mean-square amplitude,
\begin{equation}
  \rho_i^{(ab)} = \frac{\lVert\mathbf{r}_i^{(ab)}\rVert_2}{\sqrt{D}},
  \label{eq:rho}
\end{equation}
which is also the noise-free signal available to any detector that operates on mean-subtracted spectra.
To compare planets with different feature amplitudes we also form the relative residual $\mathbf{q}_i^{(ab)} = \mathbf{r}_i^{(ab)}/\bar{h}_i^{(ab)}$, with $\bar{h}_i^{(ab)} = \tfrac{1}{2}\bigl[h(\mathbf{y}_i^{(a)})+h(\mathbf{y}_i^{(b)})\bigr]$.

\paragraph{Principal components of the residuals.}
For each pair we stack the residuals into $\mathbf{Q}^{(ab)}\in\mathbb{R}^{N'\times D}$ (and analogously the absolute residuals), centre the columns, and compute the singular value decomposition $\mathbf{Q}^{(ab)}-\mathbf{1}\bar{\mathbf{q}}^{\top}=\mathbf{U}\mathbf{S}\mathbf{V}^{\top}$.
The rows $\mathbf{v}_k^{\top}$ of $\mathbf{V}^{\top}$ give the spectral shapes of the dominant inter-code differences, each signed so that its largest-magnitude loading is positive.
The scores $u_{ik}=(\mathbf{q}_i-\bar{\mathbf{q}})^{\top}\mathbf{v}_k$ give their per-planet amplitudes, and the explained-variance ratios are $S_{kk}^{2}/\sum_{k'}S_{k'k'}^{2}$.

\paragraph{Parameter dependence.}
We quantify how deterministically each component depends on the planet parameters by regressing its score on $\bm{\theta}$ with a random-forest regressor $\mathcal{R}_k:\Theta\rightarrow\mathbb{R}$ and reporting the out-of-sample coefficient of determination
\begin{equation}
  R^{2}_k = 1-\frac{\sum_{i\in\mathcal{H}}\bigl(u_{ik}-\mathcal{R}_k(\bm{\theta}_i)\bigr)^{2}}{\sum_{i\in\mathcal{H}}\bigl(u_{ik}-\bar{u}_k\bigr)^{2}},
  \label{eq:r2}
\end{equation}
on a held-out subset $\mathcal{H}$.
A value $R^2_k\rightarrow1$ indicates that component $k$ is a deterministic function of $\bm{\theta}$ and could, in principle, be emulated and corrected.
We additionally report Spearman rank correlations between the scores, the offset, and the parameters together with two derived quantities: the surface gravity proxy $\log_{10}g=\log_{10}(M_p/R_p^{2})$ in Jovian units, and $\log_{10}(T R_p/M_p)$, which is proportional to $\log_{10}(H/R_p)$ at fixed mean molecular weight, where $H$ is the atmospheric scale height.

\subsection{One-class detection with autoencoders (Q2)}
\label{sec:ae}

\paragraph{Architecture.}
For each code $c$ we train a fully connected autoencoder $f_c = g_{\bm{\psi}}\circ e_{\bm{\phi}}$ on standardised spectra from that code only.
Let $(d_1,\dots,d_K)$ be the hidden widths and $L$ the latent dimension.
With $\mathbf{a}_0=\tilde{\mathbf{y}}$, the encoder computes
\begin{align}
  \mathbf{a}_{\ell} &= \varphi\bigl(\mathbf{W}_{\ell}\mathbf{a}_{\ell-1}+\mathbf{b}_{\ell}\bigr),\qquad \ell=1,\dots,K, \label{eq:enc}\\
  \mathbf{z} &= \mathbf{W}_{K+1}\mathbf{a}_{K}+\mathbf{b}_{K+1}\in\mathbb{R}^{L}, \label{eq:latent}
\end{align}
with $\mathbf{W}_{\ell}\in\mathbb{R}^{d_{\ell}\times d_{\ell-1}}$, $d_0=D$, and $d_{K+1}=L$.
The decoder mirrors the encoder: with $\mathbf{a}'_0=\mathbf{z}$ and widths $(d_K,\dots,d_1)$,
\begin{align}
  \mathbf{a}'_{\ell} &= \varphi\bigl(\mathbf{W}'_{\ell}\mathbf{a}'_{\ell-1}+\mathbf{b}'_{\ell}\bigr),\qquad \ell=1,\dots,K, \label{eq:dec}\\
  \hat{\tilde{\mathbf{y}}} &= \mathbf{W}'_{K+1}\mathbf{a}'_{K}+\mathbf{b}'_{K+1}\in\mathbb{R}^{D}. \label{eq:out}
\end{align}
The latent and output layers are linear because standardised spectra are unbounded.
The encoder and decoder do not share weights, and the network has no normalisation, dropout, or skip connections.
Unless stated otherwise, $\varphi$ is the Gaussian error linear unit $\varphi(x)=x\,\Phi(x)$~\cite{TODO_hendrycks_gelu}, with $\Phi$ the standard normal cumulative distribution function.
We use two widths:
\begin{equation}
  \text{AE-S: } (d_1,\dots,d_4)=(128,64,32,16),\qquad
  \text{AE-W: } (d_1,\dots,d_5)=(256,128,64,32,16),
  \label{eq:widths}
\end{equation}
both with $L=8$.
The number of trainable parameters is
\begin{equation}
  P = 2\sum_{\ell=1}^{K+1}d_{\ell-1}d_{\ell} + \sum_{\ell=1}^{K+1}d_{\ell} + \sum_{\ell=0}^{K}d_{\ell},
  \label{eq:params}
\end{equation}
which gives $P=35{,}612$ for AE-S and $P=114{,}972$ for AE-W (Table~\ref{tab:arch}).
AE-W adds one 256-unit layer at each end of AE-S and has $3.2\times$ as many parameters.

\begin{table}[t]
  \centering
  \caption{Autoencoder variants. All variants are symmetric, use linear latent and output layers, and differ only in the hidden widths, the latent dimension $L$, and the activation $\varphi$. MACs: multiply--accumulate operations per forward pass.}
  \label{tab:arch}
  \begin{tabular}{lllrr}
    \toprule
    Variant & Encoder widths & $L$ & Parameters & MACs \\
    \midrule
    AE-S          & $52\!\rightarrow\!128\!\rightarrow\!64\!\rightarrow\!32\!\rightarrow\!16$                  & 8      & 35{,}612  & 35{,}072 \\
    AE-S$(L)$     & as AE-S                                                                                    & 4--16  & $35{,}348+33L$ & -- \\
    AE-W          & $52\!\rightarrow\!256\!\rightarrow\!128\!\rightarrow\!64\!\rightarrow\!32\!\rightarrow\!16$ & 8      & 114{,}972 & 113{,}920 \\
    \bottomrule
  \end{tabular}
\end{table}

\paragraph{Reconstruction objective.}
Given a mini-batch $\mathcal{B}$, the autoencoder minimises the mean squared reconstruction error in standardised units,
\begin{equation}
  \mathcal{L}_{\mathrm{rec}}(\bm{\phi},\bm{\psi}) = \frac{1}{|\mathcal{B}|\,D}\sum_{i\in\mathcal{B}}\bigl\lVert\tilde{\mathbf{y}}_i - f_c(\tilde{\mathbf{y}}_i)\bigr\rVert_2^{2}.
  \label{eq:lrec}
\end{equation}
Parameters $\bm{\omega}=(\bm{\phi},\bm{\psi})$ are updated with AdamW~\cite{TODO_loshchilov_adamw}, which applies decoupled weight decay $\lambda$,
\begin{equation}
  \bm{\omega}_{t+1} = (1-\eta_t\lambda)\,\bm{\omega}_t - \eta_t\,\frac{\hat{\mathbf{m}}_t}{\sqrt{\hat{\mathbf{v}}_t}+\epsilon},
  \label{eq:adamw}
\end{equation}
where $\hat{\mathbf{m}}_t$ and $\hat{\mathbf{v}}_t$ are the bias-corrected first and second moment estimates of $\nabla_{\bm{\omega}}\mathcal{L}_{\mathrm{rec}}$.

\paragraph{Denoising objective.}
To test whether robustness to noise can be learned, we also train denoising autoencoders~\cite{TODO_vincent_denoising_ae} with the AE-W architecture.
The input is the standardised noisy spectrum.
The target is the clean spectrum scaled by the \emph{noisy} spectrum's statistics, so that the output can be mapped back to transit depth from observable quantities only:
\begin{equation}
  \mathcal{L}_{\mathrm{den}} = \frac{1}{|\mathcal{B}|\,D}\sum_{i\in\mathcal{B}}\Bigl\lVert \frac{\mathbf{y}_i - m(\mathbf{y}_i^{\sigma})\mathbf{1}}{s(\mathbf{y}_i^{\sigma})} - f_c\bigl(\tilde{\mathbf{y}}_i^{\sigma}\bigr)\Bigr\rVert_2^{2}.
  \label{eq:lden}
\end{equation}
A new noise realisation $\bm{\epsilon}$ is drawn for every training mini-batch.

\paragraph{Anomaly scores.}
An autoencoder trained on code $c$ assigns a spectrum $\mathbf{y}$ the reconstruction score
\begin{equation}
  e_c(\mathbf{y}) = \frac{1}{D}\bigl\lVert\tilde{\mathbf{y}} - f_c(\tilde{\mathbf{y}})\bigr\rVert_2^{2}.
  \label{eq:score}
\end{equation}
Spectra from code $c$ are treated as nominal and spectra from $b\neq c$ as anomalous.
Decision thresholds $\tau_c^{(q)}$ are the $q$-th percentiles, $q\in\{95,99\}$, of $e_c$ on the nominal validation spectra, so no anomalous data are used to set them.

For noisy spectra with known per-bin uncertainty $\sigma$, $e_c$ conflates reconstruction error with noise.
We therefore also map the reconstruction back to transit depth and score it against the error bars:
\begin{equation}
  \hat{\mathbf{y}} = s(\mathbf{y}^{\sigma})\,f_c\bigl(\tilde{\mathbf{y}}^{\sigma}\bigr) + m(\mathbf{y}^{\sigma})\mathbf{1},\qquad
  \chi^{2}_{\nu,c}(\mathbf{y}^{\sigma}) = \frac{1}{D}\sum_{j=1}^{D}\Bigl(\frac{y^{\sigma}_j-\hat{y}_j}{\sigma}\Bigr)^{2}.
  \label{eq:chi2}
\end{equation}
We compare two thresholds at a target false-positive rate $\alpha$.
The first is the theoretical threshold $\tau^{\mathrm{th}} = F^{-1}_{\chi^{2}_{D}}(1-\alpha)/D$, where $F_{\chi^{2}_{D}}$ is the $\chi^{2}$ cumulative distribution function with $D$ degrees of freedom.
The second is an empirical threshold $\tau^{\mathrm{emp}}_c$, defined as the $(1-\alpha)$ quantile of $\chi^{2}_{\nu,c}$ on noisy nominal validation spectra, which absorbs the autoencoder's own reconstruction error.

\paragraph{Oracle reference.}
To separate the limits imposed by the data from those of the model, we define an oracle that scores a noisy spectrum against the true noise-free spectrum of code $c$ for the same planet:
\begin{equation}
  \chi^{2}_{\nu,c}{}^{\star}(\mathbf{y}_i^{\sigma}) = \frac{1}{D}\sum_{j=1}^{D}\Bigl(\frac{y^{\sigma}_{ij}-y^{(c)}_{ij}}{\sigma}\Bigr)^{2}.
  \label{eq:oracle}
\end{equation}
The oracle requires the true $\bm{\theta}_i$ and is not realisable in practice.
It bounds the performance of any single-spectrum goodness-of-fit test against code $c$.

\subsection{Supervised code classification (Q2)}
\label{sec:rf}

As a discriminative counterpart to the one-class detectors, we train a three-class random forest~\cite{TODO_breiman_random_forests} $\mathcal{G}:\mathbb{R}^{D}\rightarrow\Delta^{2}$ that maps a standardised (noisy) spectrum to class probabilities $\mathbf{p}=(p_{\mathrm{PSN}},p_{\mathrm{pRT}},p_{\mathrm{TRX}})$ over the three codes.
Each training planet contributes one spectrum per code, so the classifier sees every planet under every label and must rely on code-specific differences rather than planet properties.
To compare with the one-class detectors, for a pair $(a,b)$ we score test spectra of codes $a$ and $b$ with
\begin{equation}
  \pi_{ab}(\mathbf{y}) = \frac{p_b(\mathbf{y})}{p_a(\mathbf{y})+p_b(\mathbf{y})},
  \label{eq:pairscore}
\end{equation}
which plays the role of an anomaly score for nominal code $a$.

\subsection{Self- and cross-code retrieval (Q3)}
\label{sec:retrieval}

We emulate a data-driven retrieval with gradient-boosted regression trees~\cite{TODO_chen_xgboost}.
The input feature vector retains the absolute depth information removed by standardisation,
\begin{equation}
  \mathbf{x}(\mathbf{y}) = \bigl[\,\tilde{\mathbf{y}}^{\top},\ \log_{10}m(\mathbf{y}),\ \log_{10}s(\mathbf{y})\,\bigr]^{\top}\in\mathbb{R}^{D+2}.
  \label{eq:xgbfeat}
\end{equation}
For each training code $a$ and each parameter $\theta_k$, $k=1,\dots,9$, we fit an independent regressor $\hat{\theta}^{(a)}_k=G^{(a)}_k\bigl(\mathbf{x}(\mathbf{y}^{(a)})\bigr)$, which gives 27 regressors in total.
Each regressor is an additive ensemble of regression trees fitted by second-order boosting on the squared error.
Every regressor is then applied to the test spectra of all three codes for the same planets.
We call the case $b=a$ \emph{self-retrieval} and the case $b\neq a$ \emph{cross-retrieval}.
Because the planets are identical, any difference between self- and cross-retrieval errors is caused solely by $\bm{\Delta}^{(ab)}$.

For each triple $(a,b,k)$ with errors $\varepsilon_i=\hat{\theta}^{(a)}_{k}(\mathbf{y}_i^{(b)})-\theta_{ik}$ on the test set, we report the coefficient of determination $R^{2}$ (Eq.~\ref{eq:r2}), the median bias $\beta=\operatorname{median}_i\varepsilon_i$, and the half-width
\begin{equation}
  w = \tfrac{1}{2}\bigl[P_{84}(\varepsilon)-P_{16}(\varepsilon)\bigr],
  \label{eq:halfwidth}
\end{equation}
where $P_q$ is the $q$-th percentile.
We also report $\beta$ as a percentage of the parameter's range over the sample.
The ratio $|\beta|/w$ measures how large the code-induced bias is relative to the single-planet scatter.

\subsection{Diagnostics of reconstruction difficulty}
\label{sec:diag}

To identify which planets limit the autoencoders, we rank each code's held-out spectra by $e_c$ and label the top 5\% as poorly reconstructed.
For every parameter and derived quantity $\xi$, we compute the probability that a poorly reconstructed planet has a larger value of $\xi$ than another planet,
\begin{equation}
  A_{\xi} = \Pr\bigl(\xi_i > \xi_{i'} \mid i\in\mathcal{P},\ i'\notin\mathcal{P}\bigr),
  \label{eq:auc_param}
\end{equation}
where $\mathcal{P}$ is the poorly reconstructed set.
$A_\xi=0.5$ indicates no association.
We measure agreement across codes by the overlap of the sets $\mathcal{P}$ and by the Spearman correlation of per-planet scores.
An analogous breakdown of classifier accuracy is computed in quantile bins of $\bm{\theta}$, $h$, and the mean pairwise residual amplitude $\bar{\rho}_i=\tfrac{1}{3}\sum_{(a,b)}\rho_i^{(ab)}$.

% =====================================================================
\section{Experimental Setup}
\label{sec:setup}

\subsection{Data partitions}
\label{sec:splits}

We use two planet sets.
For the autoencoder design study (Section~\ref{sec:design}), the feature-height selection is applied to POSEIDON spectra only, which retains 44{,}026 planets.
For every cross-code experiment, a planet must pass the selection in all three codes, which retains 43{,}885 planets.
In both cases a single random permutation (seed 42) assigns 70\%/15\%/15\% of planets to training, validation, and test sets (30{,}818/6{,}604/6{,}604 and 30{,}719/6{,}583/6{,}583, respectively).
In the cross-code set the partition is defined on planets, so a test planet is absent from training in every code.
The validation set is used for early stopping, model selection, and threshold calibration.
The test set is used only for the reported metrics.

\subsection{Autoencoder training protocol}
\label{sec:aetrain}

Table~\ref{tab:hparams} lists the training hyperparameters.
We train for at most 1000 epochs with mini-batches of 256 spectra and use one of two learning-rate schedules.
The \emph{plateau} schedule halves $\eta$ when the validation loss has not improved for 7 epochs, with no lower bound on $\eta$.
The \emph{cosine} schedule anneals $\eta$ once per epoch $t$ as
\begin{equation}
  \eta_t = \eta_{\min} + \tfrac{1}{2}\bigl(\eta_0-\eta_{\min}\bigr)\bigl(1+\cos(\pi t/E)\bigr),
  \label{eq:cosine}
\end{equation}
with $\eta_0=10^{-3}$, $\eta_{\min}=10^{-6}$, and $E=1000$.
Early stopping tracks a reference loss $\mathcal{L}^{\mathrm{ref}}$ and resets its counter when the validation loss satisfies $\mathcal{L}_{\mathrm{val}}<(1-\delta)\,\mathcal{L}^{\mathrm{ref}}$, with relative tolerance $\delta=10^{-3}$.
Training stops after $P_{\mathrm{es}}$ consecutive epochs without such an improvement.
We use $P_{\mathrm{es}}=20$ with the plateau schedule and $P_{\mathrm{es}}=50$ with the cosine schedule.
The weights with the lowest validation loss are restored at the end of training.
Weights and biases use the default initialisation, drawn from $\mathcal{U}\bigl(-d_{\ell-1}^{-1/2},\,d_{\ell-1}^{-1/2}\bigr)$.
The initialisation seed also determines the mini-batch order.

\begin{table}[t]
  \centering
  \caption{Autoencoder training hyperparameters. Values marked $\dagger$ are framework defaults.}
  \label{tab:hparams}
  \begin{tabular}{ll}
    \toprule
    Hyperparameter & Value \\
    \midrule
    Loss                          & $\mathcal{L}_{\mathrm{rec}}$ (Eq.~\ref{eq:lrec}) or $\mathcal{L}_{\mathrm{den}}$ (Eq.~\ref{eq:lden}) \\
    Optimiser                     & AdamW, $\eta_0=10^{-3}$, $\lambda=10^{-5}$ \\
    $(\beta_1,\beta_2,\epsilon)$  & $(0.9,\,0.999,\,10^{-8})^{\dagger}$ [TODO: VERIFY framework version] \\
    Mini-batch size               & 256 \\
    Maximum epochs                & 1000 \\
    Plateau schedule              & factor 0.5, patience 7 epochs, no minimum $\eta$ \\
    Cosine schedule               & $\eta_{\min}=10^{-6}$, $E=1000$ \\
    Early stopping                & $\delta=10^{-3}$ (relative); $P_{\mathrm{es}}=20$ (plateau) or 50 (cosine) \\
    Activation $\varphi$          & GELU (exact form) \\
    Latent dimension $L$          & 8 \\
    \bottomrule
  \end{tabular}
\end{table}

\subsection{Autoencoder design study}
\label{sec:design}

The autoencoders in this work serve as measuring instruments.
A detector whose own reconstruction error exceeds the inter-code signal cannot measure that signal.
We therefore select a configuration that minimises the reconstruction error before using it for detection.
All design choices use the POSEIDON-only planet set and are made on validation MSE.
Selecting on test MSE would give the same configuration at every stage.
\begin{enumerate}
  \item \emph{Linear baseline.} Principal component analysis fitted on the training spectra, with $k\in\{4,\dots,16\}$ components, provides the reconstruction error of the optimal linear encoder of each size.
  \item \emph{Latent dimension.} AE-S with $L\in\{4,\dots,16\}$, plateau schedule, one seed (42). We select the smallest $L$ at which validation MSE reaches its plateau.
  \item \emph{Activation function.} AE-S with $L=8$ and $\varphi\in\{$ReLU, leaky ReLU, ELU, SELU, GELU, SiLU, Mish, tanh, softplus$\}$, plateau schedule, three seeds (42, 43, 44).
  \item \emph{Width and schedule.} A $2\times2$ design over \{AE-S, AE-W\} $\times$ \{plateau with $P_{\mathrm{es}}=20$, cosine with $P_{\mathrm{es}}=50$\}, seeds 42 and 43, with a third seed (44) for AE-W with the cosine schedule. The schedule and the early-stopping patience vary jointly in this design, so their individual contributions are not separated.
\end{enumerate}
The selected configuration, AE-W with GELU, $L=8$, and the cosine schedule, is used for all cross-code experiments.

\subsection{Detection experiments}
\label{sec:detsetup}

\paragraph{Noise-free detection.}
One AE-W model is trained per code on the cross-code planet set.
For POSEIDON and petitRADTRANS we train three initialisation seeds (42, 43, 44); for TauREx, one seed (42) [TODO: VERIFY whether to equalise seeds before submission].
For each code, the run with the lowest own-code validation loss is used.
We report the area under the receiver operating characteristic curve (AUC) for each of the six ordered pairs $(c,b)$, computed with the Mann--Whitney statistic
\begin{equation}
  \mathrm{AUC}_{cb} = \frac{1}{n_0 n_1}\sum_{i=1}^{n_0}\sum_{i'=1}^{n_1}\mathbb{1}\bigl[e_c(\mathbf{y}^{(b)}_{i'}) > e_c(\mathbf{y}^{(c)}_{i})\bigr],
  \label{eq:auc}
\end{equation}
together with the fractions of test spectra above $\tau_c^{(95)}$ and $\tau_c^{(99)}$.

\paragraph{Detection under noise.}
The noise-free models are applied without retraining to test spectra corrupted according to Eq.~\eqref{eq:noise} with $\sigma\in\{2,5,10,20,30,50,100\}\,\mathrm{ppm}$.
AUCs are averaged over two independent noise realisations per $\sigma$.
For the $\chi^{2}$ score we use $\alpha=0.01$, which gives $\tau^{\mathrm{th}}=1.512$.
We estimate this value from $2\times10^{6}$ samples of $\chi^{2}_{52}$.
The empirical threshold $\tau^{\mathrm{emp}}_c$ is computed on validation spectra with independent noise realisations at the same $\sigma$.
We report the true-positive rate at each threshold, the realised false-positive rate on nominal test spectra, and the oracle of Eq.~\eqref{eq:oracle}.

\paragraph{Denoising models.}
We train 12 denoising AE-W models, one for each code and each $\sigma\in\{10,20,30,50\}\,\mathrm{ppm}$, with seed 42 and otherwise the same protocol.
Validation and test spectra use one fixed noise realisation each.
These models are evaluated with the same $\chi^{2}$ protocol at the noise level used in training.

\paragraph{Random-forest classifier.}
The classifier uses 300 trees, at least 2 samples per leaf, and $\lfloor\sqrt{D}\rfloor=7$ candidate features per split, with Gini impurity and bootstrap sampling as defaults.
For each $\sigma\in\{0,10,20,30,50\}\,\mathrm{ppm}$ we add one independent noise realisation to every training and test spectrum, standardise, and train and test at the same $\sigma$.
This gives $3\times30{,}719=92{,}157$ training spectra and $3\times6{,}583=19{,}749$ test spectra.
We report accuracy, per-class recall, macro-averaged one-versus-rest AUC, and the mean pairwise AUC under the score in Eq.~\eqref{eq:pairscore}.

\subsection{Retrieval experiments}
\label{sec:retsetup}

The XGBoost~\cite{TODO_chen_xgboost} regressors use learning rate 0.1, maximum depth 8, row and column subsampling 0.8, minimum child weight 2, and histogram-based split finding.
Each is trained for up to 3000 boosting rounds, with early stopping after 100 rounds without improvement in validation RMSE.
Models are trained on noise-free spectra.
To measure how the self/cross gap scales with data volume, we retrain on random subsets containing 5\%, 10\%, 25\%, and 50\% of the training planets.
To measure the interaction between code systematics and noise, the noise-free models are applied without retraining to test spectra with $\sigma\in\{10,20,30,50\}\,\mathrm{ppm}$.
Because these models never see noise during training, their noisy-test results are a lower bound on what a noise-aware retrieval would achieve.

\subsection{Residual analysis settings}
\label{sec:ressetup}

The paired decomposition uses all 43{,}885 planets of the cross-code set without noise.
For each of the three unordered pairs we compute the first four principal components of the absolute and relative residuals.
We regress each score on $\bm{\theta}$ with a random forest of 100 trees, at least 5 samples per leaf, and half of the features per split, and report $R^2$ on a random 30\% held-out subset.

\subsection{Implementation}
\label{sec:impl}

Neural models are implemented in PyTorch~\cite{TODO_paszke_pytorch}, the random forests and metrics in scikit-learn~\cite{TODO_pedregosa_sklearn}, and the boosted trees in XGBoost~\cite{TODO_chen_xgboost} [TODO: VERIFY library versions].
All models are trained on CPU [TODO: VERIFY processor model and memory], with each autoencoder run restricted to a single thread and three to four runs executed in parallel.
One 1000-epoch run takes about 650\,s for AE-S and 920--1130\,s for AE-W.
Code and trained models are available at [TODO: VERIFY repository URL / archive DOI].

% =====================================================================
% FIGURE 1 PLACEHOLDER
\begin{figure*}[t]
  \centering
  % \includegraphics[width=\textwidth]{figures/fig1_pipeline.pdf}
  \caption{Overview of the study.
  (a) Each parameter vector $\bm{\theta}$ is passed through three radiative-transfer codes and binned to the same 52 Ariel bins, giving one spectrum per code for every planet.
  (b) Preprocessing: feature-height selection, optional white noise, and per-spectrum standardisation.
  (c) Q1: paired residuals are split into an offset and a shape term, decomposed by PCA, and regressed on $\bm{\theta}$.
  (d) Q2: one-class autoencoders, one per code, score spectra by reconstruction error or $\chi^{2}_{\nu}$ and are compared with a three-class random forest and a noise-limited oracle.
  (e) Q3: XGBoost regressors trained on one code retrieve $\bm{\theta}$ from all three codes, giving a $3\times3$ self/cross matrix.}
  \label{fig:pipeline}
\end{figure*}
